\section{A Parish Measured in Miles}

Saint Anthony's history is easiest to misunderstand when it is reduced to one founding date. A Catholic directory first places Willits within the mission field of Saint Mary's in Ukiah by 1906. Another names the Willits mission for Saint Anthony by 1910. An archival photograph record establishes that a Saint Anthony church stood in town by April 1921. In 1923 Willits Catholics petitioned for a priest of their own, an Irish Capuchin was appointed, and Sunday Mass became regular. The next annual directory available to this study lists Saint Anthony separately with its own rector. Later directories repeatedly attach ``(1923)'' to the parish name.

Each date answers a different question. The earliest checked directory appearance marks organized mission service. The patronal name marks a more definite local identity. The photograph marks a building. The petition and appointment mark resident pastoral care. The separate listing marks institutional independence in the public directory. None of them, by itself, supplies the missing canonical erection decree.

\historythesis{Saint Anthony became a parish by contraction of distance. Willits first occupied the northern edge of a large Ukiah mission circuit. A church gave that intermittent ministry a local center; the 1923 appointment of a resident Irish Capuchin changed Willits from an outlying stop into a base. From that base clergy served another wide northern field, including Covelo and stations that Catholic directories classified by nearby reservations. In 1962 the parish passed into the new Diocese of Santa Rosa; when its listed staffing changed from Capuchin to diocesan clergy around 1968--69, it retained Covelo under the name Our Lady Queen of Peace. The present English- and Spanish-language parish is therefore neither a survival untouched since one foundation day nor a sequence of disconnected institutions. It is a century-long rebalancing of center and periphery, local worship and missionary distance.}

The mileage matters. Ukiah lies south of Willits along the principal inland corridor. Covelo lies far to the northeast, beyond a direct urban-parish pattern. Earlier directories grouped still more places, rancherias, reservations, state institutions, and small settlements beneath one priest's entry. Their terse typography made a pastoral system legible to Catholic administrators, but it did not record the experience of people receiving, resisting, seeking, or sustaining ministry. The same list that proves geographical reach can hide how irregular the journey was and whose labor made a station possible.

This account therefore treats the directory as a map, not a complete parish memory. It reads that map beside Capuchin archive descriptions, official diocesan records, a tribal nation's own public history, civic records, a local newspaper, and the parish's current pages. The result is a history with strong institutional coordinates but conspicuous human absences. It can identify when Willits moved in the Catholic administrative landscape. It cannot recover every household, Indigenous encounter, sacramental journey, women's organization, building campaign, or pastoral decision from the surviving public record.

That boundary is part of the history. Saint Anthony endured not because the record is continuous, but because successive communities kept making a local Catholic center useful across discontinuities: mission to parish, Ukiah to Willits, archdiocese to diocese, Capuchin to diocesan clergy, unnamed outstation to Queen of Peace, and predominantly English public records to an explicitly bilingual present.
